% Tabla retirada de la justificación (secciones/10-justificacion-entregable.tex,
% sección Viabilidad) el 20 de octubre de 2026, a pedido del equipo.
% Es la misma tabla que sigue en secciones/07-viabilidad.tex (tab:comparativa),
% con otra etiqueta para que ambas puedan convivir en un documento.
% Para recuperarla: % Tabla retirada de la justificación (secciones/10-justificacion-entregable.tex,
% sección Viabilidad) el 20 de octubre de 2026, a pedido del equipo.
% Es la misma tabla que sigue en secciones/07-viabilidad.tex (tab:comparativa),
% con otra etiqueta para que ambas puedan convivir en un documento.
% Para recuperarla: % Tabla retirada de la justificación (secciones/10-justificacion-entregable.tex,
% sección Viabilidad) el 20 de octubre de 2026, a pedido del equipo.
% Es la misma tabla que sigue en secciones/07-viabilidad.tex (tab:comparativa),
% con otra etiqueta para que ambas puedan convivir en un documento.
% Para recuperarla: % Tabla retirada de la justificación (secciones/10-justificacion-entregable.tex,
% sección Viabilidad) el 20 de octubre de 2026, a pedido del equipo.
% Es la misma tabla que sigue en secciones/07-viabilidad.tex (tab:comparativa),
% con otra etiqueta para que ambas puedan convivir en un documento.
% Para recuperarla: \input{tablas/comparativa-tecnologias} y citarla con
% Tabla~\ref{tab:just-comparativa} (antes se citaba en Impacto Económico:
% «el código QR requiere menos infraestructura (Tabla~\ref{...})»).
\begin{table}[tbp]
	\caption{Comparación de Tecnologías de Validación de Acceso}
	\label{tab:just-comparativa}
	\begin{tabular}{lccc}
		\toprule
		Criterio                            & Código QR                & NFC     & Reconocimiento facial \\
		\midrule
		Costo de infraestructura            & Bajo                     & Medio   & Alto                  \\
		Medio que porta el estudiante       & Celular                  & Tarjeta & Ninguno               \\
		Complejidad de cómputo en el lector & Baja                     & Baja    & Alta                  \\
		Trata datos biométricos             & No                       & No      & Sí                    \\
		Riesgo de copia del identificador   & Medio\textsuperscript{a} & Bajo    & Bajo                  \\
		\bottomrule
	\end{tabular}
	\tablenote{Elaboración propia con base en
		\textcite{harnaningrum2024comparison} y \textcite{rukhiran2023user}. El costo
		y la complejidad son una valoración cualitativa del equipo.
		\textsuperscript{a}El riesgo corresponde a un QR estático; con un código
		dinámico firmado con HMAC se reduce considerablemente
		\parencite{subpratatsavee2015hmac}.}
\end{table}
 y citarla con
% Tabla~\ref{tab:just-comparativa} (antes se citaba en Impacto Económico:
% «el código QR requiere menos infraestructura (Tabla~\ref{...})»).
\begin{table}[tbp]
	\caption{Comparación de Tecnologías de Validación de Acceso}
	\label{tab:just-comparativa}
	\begin{tabular}{lccc}
		\toprule
		Criterio                            & Código QR                & NFC     & Reconocimiento facial \\
		\midrule
		Costo de infraestructura            & Bajo                     & Medio   & Alto                  \\
		Medio que porta el estudiante       & Celular                  & Tarjeta & Ninguno               \\
		Complejidad de cómputo en el lector & Baja                     & Baja    & Alta                  \\
		Trata datos biométricos             & No                       & No      & Sí                    \\
		Riesgo de copia del identificador   & Medio\textsuperscript{a} & Bajo    & Bajo                  \\
		\bottomrule
	\end{tabular}
	\tablenote{Elaboración propia con base en
		\textcite{harnaningrum2024comparison} y \textcite{rukhiran2023user}. El costo
		y la complejidad son una valoración cualitativa del equipo.
		\textsuperscript{a}El riesgo corresponde a un QR estático; con un código
		dinámico firmado con HMAC se reduce considerablemente
		\parencite{subpratatsavee2015hmac}.}
\end{table}
 y citarla con
% Tabla~\ref{tab:just-comparativa} (antes se citaba en Impacto Económico:
% «el código QR requiere menos infraestructura (Tabla~\ref{...})»).
\begin{table}[tbp]
	\caption{Comparación de Tecnologías de Validación de Acceso}
	\label{tab:just-comparativa}
	\begin{tabular}{lccc}
		\toprule
		Criterio                            & Código QR                & NFC     & Reconocimiento facial \\
		\midrule
		Costo de infraestructura            & Bajo                     & Medio   & Alto                  \\
		Medio que porta el estudiante       & Celular                  & Tarjeta & Ninguno               \\
		Complejidad de cómputo en el lector & Baja                     & Baja    & Alta                  \\
		Trata datos biométricos             & No                       & No      & Sí                    \\
		Riesgo de copia del identificador   & Medio\textsuperscript{a} & Bajo    & Bajo                  \\
		\bottomrule
	\end{tabular}
	\tablenote{Elaboración propia con base en
		\textcite{harnaningrum2024comparison} y \textcite{rukhiran2023user}. El costo
		y la complejidad son una valoración cualitativa del equipo.
		\textsuperscript{a}El riesgo corresponde a un QR estático; con un código
		dinámico firmado con HMAC se reduce considerablemente
		\parencite{subpratatsavee2015hmac}.}
\end{table}
 y citarla con
% Tabla~\ref{tab:just-comparativa} (antes se citaba en Impacto Económico:
% «el código QR requiere menos infraestructura (Tabla~\ref{...})»).
\begin{table}[tbp]
	\caption{Comparación de Tecnologías de Validación de Acceso}
	\label{tab:just-comparativa}
	\begin{tabular}{lccc}
		\toprule
		Criterio                            & Código QR                & NFC     & Reconocimiento facial \\
		\midrule
		Costo de infraestructura            & Bajo                     & Medio   & Alto                  \\
		Medio que porta el estudiante       & Celular                  & Tarjeta & Ninguno               \\
		Complejidad de cómputo en el lector & Baja                     & Baja    & Alta                  \\
		Trata datos biométricos             & No                       & No      & Sí                    \\
		Riesgo de copia del identificador   & Medio\textsuperscript{a} & Bajo    & Bajo                  \\
		\bottomrule
	\end{tabular}
	\tablenote{Elaboración propia con base en
		\textcite{harnaningrum2024comparison} y \textcite{rukhiran2023user}. El costo
		y la complejidad son una valoración cualitativa del equipo.
		\textsuperscript{a}El riesgo corresponde a un QR estático; con un código
		dinámico firmado con HMAC se reduce considerablemente
		\parencite{subpratatsavee2015hmac}.}
\end{table}
